\documentclass[11pt]{article}
\usepackage{amsmath,amssymb,amsthm}
\usepackage[margin=1in]{geometry}
\newtheorem{theorem}{Theorem}
\newtheorem{lemma}{Lemma}
\newtheorem{corollary}{Corollary}
\newtheorem{definition}{Definition}
\title{V99: Unbalanced Singleton Cores and the Signed-Majority $X$ Dichotomy}
\author{NC0\_k-Avoid Laboratory}
\date{}
\begin{document}
\maketitle

\section{Singleton-core gates}
Let a local essential ternary gate lie in the NPN orbit of $\mathrm{AND}_3$.
Choose an output orientation so that the oriented gate $h_e$ has exactly one
satisfying local assignment $\alpha_e$. Then
\[
 h_e(x)=1 \iff x_v=\alpha_e(v)\quad\text{for every }v\in\operatorname{supp}(e).
\]

\begin{lemma}[Conflict certificate]
If two oriented singleton-core gates share an input $v$ and demand opposite
values there, then requesting both oriented outputs to be $1$ gives a missing
output word.
\end{lemma}
\begin{proof}
A realizing input would have to assign $v$ both Boolean values. Undoing the
chosen output complements transports the certificate to the original circuit.
\end{proof}

\begin{lemma}[Conflict-free switching]
If no shared input receives contradictory singleton demands, then an input
switching transforms every oriented gate into the positive conjunction of its
essential inputs.
\end{lemma}
\begin{proof}
For each active input $v$, let $r_v$ be its common demanded bit and put
$z_v=1$ iff $x_v=r_v$. Every oriented singleton gate becomes
$\bigwedge_{v\in\operatorname{supp}(e)} z_v$.
\end{proof}

\begin{theorem}[Singleton-core polynomial avoidance]
Every positive-surplus $NC^0_3$ component whose essential gates all lie in the
singleton NPN orbit admits deterministic polynomial-time range avoidance.
\end{theorem}
\begin{proof}
If a singleton-demand conflict exists, use the first lemma. Otherwise use the
second lemma and apply the deterministic polynomial-time monotone-$NC^0_3$
avoider of Kuntewar--Sarma for $m>n$, then undo the output orientations.
\end{proof}

\begin{theorem}[Strict extension beyond V98]
For every $N\ge5$ there is a connected exact-stretch singleton-core circuit
with V97 peeling parameter $\lambda=N$ that is not switching-balanced in the
V98 sense but has an immediate singleton conflict certificate.
\end{theorem}
\begin{proof}
Use cyclic supports $S_i=\{i,i+1,i+2\}\pmod N$ with positive $\mathrm{AND}_3$
and one duplicate of $S_0$ labeled
$(\neg x_0)\wedge x_1\wedge x_2$. Every input has incidence degree at least
three, so V97 does not peel. The two copies of $S_0$ demand opposite values of
$x_0$ when their oriented outputs are $1$. The V98 incidence-sign equations are
inconsistent because the positive copy forces equal switching potentials on
$x_0,x_1,x_2$, whereas the duplicate requires the sign pattern $(1,0,0)$.
\end{proof}

\section{Signed-majority simple $X$}
Let $L=2\ell\ge6$. The simple $X_L$ has cycle vertices
$v_0,\dots,v_{L-1}$, special edges
\[
 e_0=\{v_0,v_1,w\},\qquad e_3=\{v_3,v_4,w\},
\]
and every other edge $e_i=\{v_i,v_{i+1},p_i\}$ with private third vertex
$p_i$. Each edge is labeled by majority of three signed literals.

Because
\[
 \neg\operatorname{MAJ}(a,b,c)=\operatorname{MAJ}(\neg a,\neg b,\neg c),
\]
input and output complements act as vertex switching on the signed bipartite
incidence graph.

\begin{lemma}[Four switching classes]
The signed incidence graph of $X_L$ has exactly four switching classes.
They are indexed by two independent cycle parities $(B,A)\in\mathbb F_2^2$ and
have canonical representatives with all incidence signs positive except
possibly $(e_0,v_0)=B$ and $(e_0,w)=A$.
\end{lemma}
\begin{proof}
The connected incidence graph has $3L-1$ vertices and $3L$ incidences, hence
cycle rank two. Switching classes of binary edge signs are therefore indexed by
the two cycle parities. A spanning-tree switching normal form leaves exactly the
two stated cotree signs.
\end{proof}

For a positive private majority edge with cycle endpoints $a,b$ and existential
private input $p$, target $0$ is feasible except at $(a,b)=(1,1)$ and target
$1$ is feasible except at $(a,b)=(0,0)$. Hence the endpoint relations are
\[
 R_0=\begin{pmatrix}1&1\\1&0\end{pmatrix},\qquad
 R_1=\begin{pmatrix}0&1\\1&1\end{pmatrix}
\]
over the Boolean semiring.

\begin{lemma}[Path collapse]
If a private-edge target path contains equal adjacent target bits, its endpoint
relation is universal. If it is alternating and has positive even length, its
relation is
\[
 P_{01}=\begin{pmatrix}1&1\\0&1\end{pmatrix}
 \quad\text{or}\quad
 P_{10}=\begin{pmatrix}1&0\\1&1\end{pmatrix}
\]
according to its first target bit.
\end{lemma}
\begin{proof}
Boolean matrix multiplication gives
$R_0R_0=R_1R_1=J$, $JR_b=R_bJ=J$,
$R_0R_1=P_{01}$, $R_1R_0=P_{10}$, and
$P_{01}^2=P_{01}$, $P_{10}^2=P_{10}$.
\end{proof}

\begin{lemma}[Constant boundary table]
After collapsing the two even private paths, the only infeasible boundary
configurations are, in switching class $(B,A)=(0,0)$,
\[
 (P_{01},P_{10},y_0,y_3)=(P_{01},P_{10},1,0)
\]
and
\[
 (P_{10},P_{01},y_0,y_3)=(P_{10},P_{01},0,1).
\]
The three nonzero switching classes have no infeasible boundary configuration;
if either path relation is universal, every boundary configuration is feasible.
\end{lemma}
\begin{proof}
Only the five variables $v_0,v_1,v_3,v_4,w$ remain. Substitution into
\[
 \operatorname{MAJ}(v_0\oplus B,v_1,w\oplus A)=y_0,\qquad
 \operatorname{MAJ}(v_3,v_4,w)=y_3
\]
together with the two $2\times2$ endpoint relations gives a constant
$32$-assignment Boolean check, independent of $L$.
\end{proof}

\begin{theorem}[Signed-majority $X$ dichotomy]
For every even $L\ge6$, a signed-majority simple $X_L$ satisfies:
\[
\begin{array}{c|c}
\text{switching class}&\text{missing edge words}\\\hline
(0,0)&2\\
(0,1),(1,0),(1,1)&0.
\end{array}
\]
In the all-positive representative the two missing words are the two alternating
edge words. Every nonzero switching class is surjective.
\end{theorem}
\begin{proof}
Both private paths have even length. The path-collapse lemma reduces every
target to the constant boundary table. The two failures in the balanced class
are exactly the two globally alternating targets; all other targets are
realizable. The three nonzero classes have no failing boundary type. Switching
is a bijection on the input cube followed by output-coordinate complements, so
it preserves range cardinality and surjectivity.
\end{proof}

\paragraph{Boundary.}
This theorem classifies a local loose-$X$ certificate, not the full
signed-majority circuit class. The remaining essential ternary unate NPN orbit
of size $48$, represented by $x\wedge(y\vee z)$, is left open. No unrestricted
$NC^0_3$ algorithm, lower-bound transfer, novelty, or P-versus-NP claim is made.
\end{document}
